\begin{figure}[htbp]
\centering
\begin{tikzpicture}[x=0.92cm,y=0.85cm,>=stealth,
  every node/.style={font=\small}]
\foreach \y in {0,1.3,2.6} {
  \draw[->,gray] (-0.4,\y)--(7,\y) node[right] {$x$};
}
\draw[very thick,arkivblaa] (1,2.6)--(5.6,2.6);
\draw[thick,draw=arkivblaa,fill=white] (1,2.6) circle (2.7pt);
\fill[arkivblaa] (5.6,2.6) circle (2.7pt);
\node[below] at (1,2.5) {$a$};
\node[below] at (5.6,2.5) {$b$};
\node[anchor=east] at (-0.6,2.6) {$(a,b]$};
\draw[very thick,arkivblaa] (1,1.3)--(3.2,1.3);
\draw[very thick,arkivblaa] (3.2,1.3)--(5.6,1.3);
\foreach \x in {1,5.6} {
  \draw[thick,draw=arkivblaa,fill=white] (\x,1.3) circle (2.7pt);
}
\draw[thick,draw=arkivrod,fill=white] (3.2,1.3) circle (2.7pt);
\node[below] at (1,1.2) {$a$};
\node[below] at (3.2,1.2) {$c$};
\node[below] at (5.6,1.2) {$b$};
\node[anchor=east] at (-0.6,1.3) {$(a,b)\setminus\{c\}$};
\draw[very thick,arkivblaa] (1.4,0)--(5.2,0);
\foreach \x in {1.4,5.2} {
  \draw[thick,draw=arkivblaa,fill=white] (\x,0) circle (2.7pt);
}
\draw (3.3,-0.08)--(3.3,0.08);
\node[below] at (3.3,-0.1) {$x_0$};
\draw[<->,arkivrod] (3.3,0.4)--(5.2,0.4)
  node[midway,above] {$r$};
\node[below] at (1.4,-0.1) {$x_0-r$};
\node[below] at (5.2,-0.1) {$x_0+r$};
\node[anchor=east] at (-0.6,0) {$|x-x_0|<r$};
\end{tikzpicture}
\caption{Intervaller på en skematisk, dimensionsløs talakse. En fyldt prik
medtager et endepunkt, mens en åben prik udelader det. I midten antages
$a<c<b$: udeladelsen af $c$ deler området i to intervaller. Nederst vises
afstandsbetingelsen omkring $x_0$ med $r>0$.}
\label{fig:algebra-intervaller}
\end{figure}
